	\section{Introduction}

Reading is often described as "the process of gaining meaning from print" \citep[][p.~34]{raynerHowPsychologicalScience2001}. Although comprehension ultimately depends on sentence- and discourse-level integration, efficient access to individual word forms and meanings remains a foundational component of skilled reading. However, readers vary considerably in this ability \citep{perfettiReadingAbilityLexical2007, castlesHowDoesOrthographic2006, stanovichExposurePrintOrthographic1989}. Understanding what drives this variation requires identifying variables that are selectively diagnostic of the mechanisms underlying word recognition. Two such variables, base frequency (also referred to as stem frequency in the literature) and morphological family size, have proven especially informative for this purpose \citep{taftMorphologicalDecompositionReverse2004, taftRecognitionAffixedWords1979, dejongMorphologicalFamilySize2000, bertramEffectsFamilySize2000}. Base frequency, the cumulative token frequency of a base morpheme across all its surface forms, and morphological family size, the number of distinct words sharing that base, both facilitate visual word recognition, yet their effects are independent and have been argued to reflect different levels of morphological representation: base frequency has been argued to index form-based access, reflecting the cumulative ease with which a base morpheme's orthographic form is processed across its surface variants, while family size has been argued to index the richness of semantic representation, reflecting the breadth of meaning-level connections activated when a morphological base is encountered \citep{schreuderHowComplexSimplex1997}. The present study exploits this dissociation to ask how readers who differ in the precision of their orthographic and semantic knowledge draw differentially on form-based and meaning-based information during morphological processing.

Both variables are properties of a word's morphological family, and morphological structure plays a particularly important role in the English writing system. English exhibits irregular and context-sensitive grapheme–phoneme mappings, rendering phonological decoding alone an insufficient basis for fluent word recognition \citep{zieglerReadingAcquisitionDevelopmental2005, castlesHowDoesOrthographic2006}. At the same time, English preserves regularities at higher representational levels, particularly in its morphological structure; words that share a stem or affix often exhibit consistent orthographic patterns and systematic semantic relationships (e.g., \textit{heal–health} share a stem; \textit{teacher}, \textit{baker}, and \textit{driver} share the agentive suffix \textit{-er}). Skilled reading in English therefore requires sensitivity not only to phonological correspondences but also to recurring morphological structure that links form and meaning, and readers differ in how deeply and precisely that structure is represented \citep{bryantMorphologySpellingWhat2005, arciuliReadingStatisticalLearning2018, treimanStatisticalLearningSpelling2018}.

Base frequency and family size are the primary variables used to study this structural sensitivity; before characterising how readers differ in processing them, however, it is necessary to clarify what each variable indexes and how it has been theoretically interpreted. This question has proven theoretically contentious because the answer depends on the architecture one ascribes to the word recognition system. Whether skilled readers decompose complex words into constituent morphemes, access whole-word representations directly, or exploit learned statistical regularities carries different implications for what each variable measures, that is, whether it reflects the strength of a stored morpheme entry, the output of a parsing procedure, or an emergent property of a statistical learning system. As a result, what each variable measures determines what kind of reader variation should modulate its effect. The Lexical Quality Hypothesis \citep{perfettiLexicalQualityHypothesis2002, perfettiReadingAbilityLexical2007, perfettiLexicalQualityRevisited2017} provides a principled framework for characterising that variation: lexical representations vary across readers in how completely and precisely a word's orthographic and semantic constituents are specified, and it is this variation in representational quality that determines how strongly each variable influences processing. Although competing architectures disagree about the mechanism responsible, they converge in treating base frequency as an index of how precisely the orthographic form of the base morpheme is specified, and family size as an index of how richly and flexibly its meaning is represented across morphological contexts. Readers who differ in orthographic precision versus semantic breadth should therefore differ selectively in their sensitivity to each variable.


\subsection{Base Frequency, Morphological Family Size, and Competing Theoretical Accounts}

Both base frequency and morphological family size facilitate visual word recognition, and their effects are independent when appropriately controlled \citep{fordDerivationalMorphologyBase2010}. The theoretical significance of this dissociation is that it constrains what any model of morphological processing must explain: current accounts converge on the prediction that base frequency effects arise earlier in processing than family size effects, but diverge sharply in the mechanisms responsible.

Dual-route and multi-stage accounts agree that the two variables index distinct levels of processing, but disagree on the nature of those levels. In classical dual-route frameworks \citep{schreuderModelingMorphologicalProcessing1995, dejongMorphologicalFamilySize2000}, base frequency reflects the ease of accessing a stored morpheme entry via a decompositional route, while family size reflects the richness of the semantic network surrounding the base, attributed to the whole-form route. Early automatic decomposition accounts \citep{taft_lexical_1975, taftMorphologicalDecompositionReverse2004} reconceptualise this dissociation in temporal rather than architectural terms: decomposition is not one route among several but an obligatory initial parsing step that precedes semantic access. On the morpho-orthographic formulation of this view \citep{rastleMorphologicalDecompositionBased2008, longtinMorphologicalDecompositionEarly2005, morrisSemanticTransparencyMasked2007}, an intermediate level of representation mediates between sublexical letter clusters and whole-word forms, and base frequency reflects the strength of activation at this stage. Family size effects are then linked to later semantic processing, consistent with the finding that opaque relatives contribute less to the family size effect than transparent ones \citep{schreuderHowComplexSimplex1997, dejongMorphologicalFamilySize2000}. 

The discriminative learning framework \citep[NDL/LDL;][]{chuangDiscriminativeLearningLexicon} offers a more radical reinterpretation. On this account, lexical processing is modelled as learned form-to-meaning mappings driven by prediction error rather than by decompositional morpheme retrieval. Base frequency reflects the accumulated strength of the n-gram cues that spell out the base across all its surface forms, while family size reflects the type diversity of the training signal; a larger family causes shared cues to converge on a more robustly specified semantic core. Both effects arise from the same learning mechanism, rendering the dissociation between them a consequence of different statistical properties of the input rather than evidence for distinct representational levels or processing routes.

Taken together, these accounts agree that base frequency and family size exert independent effects, with base frequency arising earlier in processing than family size. Where they diverge is in the nature of the representations involved---symbolic morpheme entries, learned orthographic patterns, or weighted form-to-meaning mappings---and consequently in how they explain the source of individual differences in sensitivity to each variable. Dual-route and morpho-orthographic accounts imply that readers with more precise sublexical orthographic representations should show greater sensitivity to base frequency, since this variable indexes activation strength at the form-parsing stage, while readers with richer semantic networks should show greater sensitivity to family size, since this variable indexes the breadth of meaning-level connectivity. The discriminative learning framework leads to a similar expectation but for different reasons: individual differences in statistical learning history determine the specificity of both orthographic cue weights and semantic vectors, so readers with more extensive print exposure should show stronger and more differentiated sensitivity to both variables. However, all three accounts are consistent with the view that form-level and meaning-level sensitivity should be dissociable, and that this dissociation should be traceable to distinct processing stages---early for base frequency, later for family size.

Across these accounts, a recurring question is whether morphological decomposition is driven primarily by affix recognition or by stem activation. Affix-stripping accounts \citep{taft_lexical_1975, rastleMorphologicalDecompositionBased2008} propose that the affix is identified first and triggers decomposition, after which the stem is looked up in the lexicon. However, \citet{beyersmannEmbeddedStemPriming2016} found that high proficiency readers showed embedded stem priming effects independent of affix status: priming occurred even when the stem was followed by a non-affixal ending, suggesting that stem activation in skilled readers is not contingent on affix recognition. This finding suggests that affix-stripping alone may be insufficient to account for morphological processing in skilled readers, and raises the possibility that stem- and affix-based mechanisms operate concurrently, with their relative contributions varying as a function of individual differences in reading proficiency. The present study tests this possibility directly by examining whether Orthographic Sensitivity and Language Proficiency, two dissociable dimensions of individual variation in reading skill, differentially modulate early morpho-orthographic parsing in ways consistent with affix-based and stem-based mechanisms respectively.


\subsection{Individual Differences in Morphological Processing}

Traditional models of word recognition have typically characterised processing at the level of the skilled reader in general, abstracting over individual differences in favour of a uniform cognitive architecture. Yet readers vary considerably in how they process written words, and understanding this variation offers a window into the representations and processes underlying skilled reading. The Lexical Quality Hypothesis, introduced above, provides a principled framework for understanding this variation. Readers do not differ along a single dimension of lexical quality. Andrews and colleagues  \citep{andrewsIndividualDifferencesSkilled2012, andrewsMorphologicalPrimingStronger2013} have demonstrated that individual differences in orthographic and semantic knowledge are dissociable and predict distinct patterns of morphological processing. Readers with a semantic profile, defined by relatively higher vocabulary than spelling ability, show robust priming for transparent morphological pairs but little priming for opaque or form-only pairs, particularly for slower responses. Readers with an orthographic profile, defined by relatively higher spelling than vocabulary, show sustained priming for opaque pairs that is at least as strong as for transparent pairs. \citet{andrewsMorphologicalPrimingStronger2013} interpret this pattern as evidence that a systematic component of individual differences among skilled readers is explicable in terms of the balance between orthographic and semantic processing; readers differ not merely in overall processing efficiency but in the relative weight they assign to form-based versus meaning-based information during lexical access.

The present study builds directly on this finding. If base frequency and family size index form-based and meaning-based aspects of lexical representation respectively, then readers who differ in their balance of orthographic and semantic processing should differ predictably in their sensitivity to each variable. Moreover, if these differences have a temporal structure, with form-based processing preceding semantic processing in the recognition cascade, they should be traceable to distinct ERP components. To capture the full range of this variation continuously rather than as a categorical split, the present study derives two orthogonal individual difference dimensions, Language Proficiency and Orthographic Sensitivity, from a principal components analysis of vocabulary knowledge, spelling ability, and print exposure, following the approach of \citet{andrewsMorphologicalPrimingStronger2013}. If base frequency and family size are differentially modulated by these two dimensions, and if those modulations emerge at distinct stages in the ERP, this would constitute direct electrophysiological evidence that orthographic and semantic processing represent a genuine and temporally structured axis of variation in the human reading system.


\subsection{Neural Indices of Morphological Processing}

To track whether base frequency and family size effects emerge at distinct temporal stages, and whether individual differences in orthographic and semantic processing modulate those stages selectively, the present study uses two well-established ERP components as dependent measures---the N250, as an index of sublexical processing, and the N400, as an index of lexical-semantic integration. These components have been argued to reflect sequential stages in the cascade from orthographic decoding to meaning retrieval during visual word recognition \citep{graingerWatchingWordGo2009}.

The N250 is a negative-going deflection peaking at approximately 250 ms post-stimulus onset. \citet{holcombTimeCourseVisual2006} propose that it reflects the stage at which abstract sublexical orthographic representations, that is letters and letter clusters, are mapped onto higher-level form representations prior to full lexical access. Evidence for this interpretation comes from its graded sensitivity to orthographic overlap between prime and target in masked priming; N250 amplitude is most attenuated for full identity repetitions, intermediate for primes sharing all but one letter, and largest for fully unrelated pairs \citep{holcombTimeCourseVisual2006, holcomb_exploring_2007}. The N250 is also sensitive to sublexical phonological overlap, with \citet{graingerTimeCourseOrthographic2006} identifying separable early orthographic and later phonological subcomponents, consistent with dual-route architectures of word recognition.

The N400 is a negative-going component peaking around 400 ms, with a broad centro-parietal distribution \citep{kutasReadingSenselessSentences1980}. Across more than four decades of research, its amplitude has been shown to vary with semantic priming, cloze probability, lexical association, concreteness, and semantic richness, while remaining insensitive to physical and syntactic manipulations that leave meaning intact \citep{kutasThirtyYearsCounting2011}. For the present study, the N400 is reliably modulated in single-word and word-pair paradigms, with larger amplitudes for semantically unrelated than related targets, indexing the ease with which a word's meaning is integrated with concurrently activated semantic information: in the present design, the network of meanings associated with a stem's morphological family, rather than sentence- or discourse-level content.

Although both components are negative-going deflections occurring in adjacent time windows, they are dissociable by topography---the N250 has a more frontal distribution, the N400 a posterior one \citep{holcombTimeCourseVisual2006}---and, more decisively, by a double dissociation across priming types. Lexical status selectively modulates the N400 in that priming effects are restricted to real words with semantic representations and do not extend to pseudowords \citep{kiyonagaMaskedCrossmodalRepetition2007}, whereas the N250 is equally modulated by words and pseudowords alike. The reverse pattern holds for semantic relatedness; while associative priming reliably modulates the N400, it produces no N250 effect \citep{graingerWatchingWordGo2009}. This double dissociation---pseudoword priming engaging the N250 but not the N400, semantic priming engaging the N400 but not the N250---confirms that the two components index qualitatively distinct processing stages, and rules out the interpretation of the N250 as merely an early onset of the N400.
